\section*{PROJECT SUMMARY}
%\vspace{-.1in}

The goal of this research project is to build an infrastructure
that helps multimedia viewers interact with the content and 
change the pace of a televised program.
A fully interactive system (e.g., telepresence) would allow
the person to be immersed in a distant environment
and act upon it as if he was there, but we are far from
that capacity given the technology of the near future.
In this proposal, I present a receiver-based
architecture that supports {\em viewing interactivity}
and {\em content interactivity}. Viewing interactivity enables
a viewer to time-shift a televised program through
pause, fast-forward and instant replay operations. 
Content interactivity allows a viewer to query or to
change the content of multimedia data. 

With viewing interactivity, one can watch a program at 
whatever pace one finds comfortable. For instance, 
when one misses a crucial point while viewing a televised
lecture, one can replay the missed section and then 
continue viewing the lecture with a phase delay, or fast-forward
the lecture to get back in synch with the televised stream.
One can also record programs at multiple channels at the same
time, filter out unwanted
content and produce a customized program. 
(Using today's tape-based VCR, one cannot record
and replay a program or record more than one program 
at the same time.)
Content interactivity enables one to introduce one's 
creativity into a program.
For example, children would be able to change the attributes 
(e.g., colors, size, position and even the personality)
of the ``bunny'' and ``big bird'' 
characters in the {\it Sesame Street} program.  
With the capabilities to change both pace and content,
a class of high school students could conduct science experiments
in small groups through a virtual lab and their instructor
could keep track of their progress in real-time.   
Many other applications exist.

To enable viewing and content interactivity, 
we need a multimedia receiver that can display
more than one stream of TV, radio and the Web data,
can time-shift a stream, and can query and change
the content of a stream.  To accomplish these
goals, the multimedia receiver
needs a {\em buffer manager} that can schedule the
receiver's memory and 
disk bandwidth effectively to satisfy 
various interactive scenarios, and it needs a {\em data manager}
that can manage data from TV, radio and the Web 
for supporting efficient queries and retrieval.

Although progress has been made in buffer and
data management, many issues, especially in managing multimedia
data and managing buffer for rendering and
for querying multimedia data,
and in developing an architecture that can be adaptive
to both resource constraints and {\em Quality of
Service} (QoS) requirements, remain open.
I plan to address these issues, and to build on
work from several fields, most notably from the
image processing, operating system and database fields.
In particular, I am interested in developing the 
following technologies:

\begin{nlist}
\item {\bf 1}.
Temporal indexing structures that support
efficient stream synchronization.
\item {\bf 2}. Caching and prefetching techniques that satisfy
the real-time data requirement of the continuous media data
(i.e., video and audio) on the one hand, and minimize the
initial latency in responding to an interactive request
on the other.
\item {\bf 3}.  
QoS-adaptive buffer-management architectures that manage
memory and disk bandwidth by adapting to
the receiver's resource constraints and
the viewer's QoS requirements.
\item {\bf 4}.
Media query processing that handles heterogeneous
media queries efficiently. 
\end{nlist}

In addition, I envision constructing demonstrations 
that use these technologies.
Here is a list of some of the demonstrations I am planning
to indicate the progress of this project:
\begin{jlist}
\item 
A multimedia database that supports content-based 
queries,
\item 
A digital VCR that supports pause, fast-forward and instant 
replay, and
\item 
A multi-camera video surveillance system (e.g., for 
monitoring in a nuclear power facility).
\end{jlist}

%The multimedia database system provides the digital-VCR 
%project with an efficient storage and an effective query 
%engine that can help the viewer select content 
%before the viewer issues a fast-forward or replay command.
%The digital-VCR, which receives only one stream at a time,
%in turn, is extended by the multi-camera video surveillance project,
%which handles multiple cameras, streams and viewers simultaneously.

Finally, my education objectives include the expansion of
the curriculum in database and multimedia systems at U.C. Santa
Barbara. In particular, a new course that I am developing
will cover the latest 
trends in database and multimedia systems. 

In summary, this project will develop an infrastructure 
that provides a framework for building
a wide variety of interactive multimedia applications
in the educational, commercial, government and
military sector.

